% !PRINTED 573
\weberpage{573}

Nach §\,150, (8), (10), (13) beginnt die Entwickelung von $\tau(u)$ nach steigenden Potenzen von $u$ mit einer negativen Potenz von $u$. Wir ordnen in der Gleichung

\weq{4}{\tau(\mu u)\, N = Z}

die Funktion $Z$ nach fallenden Potenzen von $\tau$, entwickeln beide Seiten nach steigenden Potenzen von $u$ und vergleichen die Koeffizienten gleich hoher Potenzen von $u$. Dann bekommen wir für die Bestimmung der Koeffizienten von $Z$ eine Reihe linearer Gleichungen, deren jede folgende nur einen neuen dieser Koeffizienten enthält, und daraus können diese Koeffizienten successive rational berechnet werden. Beachtet man nun die Form der Entwickelungen des §\,150, so ergibt sich, daß die Koeffizienten von $Z$ rationale Funktionen von $j(\omega)$ und $\sqrt{\Delta}$ sind.

Befreit man nach dem Euklidischen Algorithmus $N$ und $Z$ von gemeinschaftlichen Faktoren, so kommt man zu den Funktionen $\Psi, \Phi$ zurück, und es ergibt sich, wenn wir unter dem Klassenkörper den Inbegriff der rationalen Funktionen von $\sqrt{\Delta}$ und $j(\omega)$ verstehen:

\begin{quote}
3. \gesperrt{Die Koeffizienten der Funktionen $\Phi(\tau)$, $\Psi(\tau)$ in (2) gehören dem Klassenkörper an.}
\end{quote}

\begin{center}
§\,154. \textbf{Der Teilungskörper.}
\end{center}

Wenn wir (durch rationale Rechnung) die Funktion $\Phi(x)$ von allen mehrfach darin vorkommenden Faktoren befreien, so bleibt eine ganze Funktion $T_\mu(x)$ von $x$ übrig, deren Koeffizienten dem Klassenkörper angehören, und die Wurzeln von $T_\mu(x)$ sind sämtliche voneinander verschiedene unter den Größen

\weq{1}{\tau\left(\frac{\nu}{\mu}\right).}

Zwei dieser Größen

\[
\tau\left(\frac{\nu}{\mu}\right), \quad \tau\left(\frac{\nu'}{\mu}\right)
\]

sind einander gleich,

% [sic] Im Fall b) druckt die Quelle auch vor +/-i*nu'/mu ein Kongruenzzeichen; die Alternativen sind in (2) auf der Folgeseite durch Kommas getrennt. Das gedruckte Zeichen bleibt erhalten.
\[
\begin{array}{rl}
\text{im Falle (1) a), wenn } & \dfrac{\nu}{\mu} \equiv \pm\dfrac{\nu'}{\mu}, \\
\text{(1) b), wenn } & \dfrac{\nu}{\mu} \equiv \pm\dfrac{\nu'}{\mu} \equiv \pm i\dfrac{\nu'}{\mu}, \\
\text{(1) c), wenn } & \dfrac{\nu}{\mu} \equiv \pm\dfrac{\nu'}{\mu}, \quad \pm\varrho\dfrac{\nu'}{\mu}, \quad \pm\varrho^2\dfrac{\nu'}{\mu},
\end{array}
\]
%CONT>>
